\documentclass[12pt,a4paper,UTF8]{ctexart}
\usepackage[a4paper,top=22mm,bottom=22mm,left=22mm,right=22mm,headheight=15pt]{geometry}
\newcommand{\Needspace}[1]{\par\begingroup\dimen0=#1\relax\dimen2=\pagegoal\advance\dimen2 by-\pagetotal\ifdim\dimen0>\dimen2\newpage\fi\endgroup}
\newcommand{\needspace}[1]{\par\begingroup\skip0=\lastskip\vskip-\skip0\dimen0=#1\relax\advance\dimen0 by\skip0\vskip\dimen0\penalty-100\vskip-\dimen0\vskip\skip0\penalty9999\endgroup}
\usepackage{amsmath,amssymb,mathtools,bm}
\usepackage{graphicx,xcolor,booktabs,longtable,array,enumitem,fancyhdr,caption,float,microtype,setspace}
\usepackage[hidelinks,unicode]{hyperref}
\usepackage[most]{tcolorbox}
\definecolor{ContestBlue}{HTML}{17365D}
\hypersetup{pdftitle={第055套},pdfauthor={}}
\setstretch{1.24}
\setlength{\parindent}{2em}
\setlength{\parskip}{0.35em}
\setlength{\emergencystretch}{2em}
\setlist{itemsep=0.35em,topsep=0.4em}
\captionsetup{font=small,labelfont=bf,justification=centering}
\setcounter{secnumdepth}{0}
\pagestyle{fancy}
\fancyhf{}
\fancyhead[L]{\small 第055套}
\fancyfoot[C]{\thepage}
\renewcommand{\headrulewidth}{0.4pt}
\ctexset{
  section={format=\Large\bfseries\color{ContestBlue},beforeskip=1.2em,afterskip=0.7em},
  subsection={format=\large\bfseries\color{ContestBlue},beforeskip=1.0em,afterskip=0.5em},
  subsubsection={format=\normalsize\bfseries,beforeskip=0.8em,afterskip=0.35em}
}
\usepackage{calc}
\newcounter{none}
\let\originalsubsection\subsection
\renewcommand{\subsection}{\needspace{6\baselineskip}\originalsubsection}
\makeatletter
\newcommand{\tightlist}{\setlength{\itemsep}{0pt}\setlength{\parskip}{0pt}}
\newcommand{\maxwidth}{\ifdim\Gin@nat@width>\linewidth\linewidth\else\Gin@nat@width\fi}
\newcommand{\maxheight}{\ifdim\Gin@nat@height>0.72\textheight 0.72\textheight\else\Gin@nat@height\fi}
\newcommand{\pandocbounded}[1]{\begingroup\setkeys{Gin}{width=\maxwidth,height=\maxheight,keepaspectratio}#1\endgroup}
\makeatother
\begin{document}
\begin{titlepage}
  \centering
  \vspace*{58mm}
  {\Huge\bfseries\color{ContestBlue} 第055套\par}
  \vfill
  \vspace*{28mm}
\end{titlepage}
\section{第 1 题}\label{ux7b2c-1-ux9898}

空间中存在匀强磁场 \(\bm B\)，磁场位于 \(x\)-\(z\) 平面内，并与 \(z\)
轴夹角为 \(\theta\)。质量为 \(m\)、电荷量为 \(q\)
的点电荷从原点出发，初速度为 \(\bm v_0=v_0\bm e_x\)。求：

（1）粒子的运动轨迹参数方程和长时间平均速度；

（2）粒子回旋运动的磁矩大小；

（3）若磁场方向以角速度 \(\omega\ll |q|B/m\) 缓慢绕 \(z\)
轴转动，并忽略感生电场，选取满足

\[
\frac{2\pi m}{|q|B}\ll\tau\ll\frac{2\pi}{\omega}
\]

的时间窗，求该时间窗内的平均速度；

（4）静态磁场中，粒子是否会在 \(t>0\) 时回到 \(z\) 轴。

\begin{figure}[H]
\centering
\includegraphics[width=0.529\linewidth,height=\textheight,keepaspectratio,alt={图1-1　斜磁场方向、初速度与螺旋运动的空间关系}]{../figures/question-01.pdf}
\caption*{图1-1　斜磁场方向、初速度与螺旋运动的空间关系}
\end{figure}

\newpage

\section{第 2 题}\label{ux7b2c-2-ux9898}

Michelson 干涉仪使用半反射板 \(G_1\)
将入射光分成互相垂直的两束光，分别经反射镜 \(M_1\)、\(M_2\)
反射后重新会合，\(G_2\) 为补偿板。\(M_2\) 固定，\(M_1\)
可沿光束方向移动。求：

（1）使用波长 \(589.3\,\mathrm{nm}\) 的钠光，移动 \(M_1\) 后中心吞吐
\(20\) 圈条纹，\(M_1\) 移动的距离；

（2）钠黄光由两条相近谱线组成，平均波长
\(\lambda_0=589.3\,\mathrm{nm}\)。条纹从最模糊到最清晰再到最模糊的一个周期内吞吐
\(985\) 圈，求双线波长差；

（3）以速度 \(v\) 匀速移动 \(M_1\)，光电信号频率为 \(f\)，求波长
\(\lambda\)，并对
\(v=6.0\times10^{-4}\,\mathrm{m/s}\)、\(f=2.54\times10^3\,\mathrm{Hz}\)
求数值；

（4）按上述速度移动反射镜，求钠黄光双线电信号的拍频。

\begin{figure}[H]
\centering
\includegraphics[width=0.546\linewidth,height=\textheight,keepaspectratio,alt={图2-1　Michelson 干涉仪的分束、往返反射和观察方向}]{../figures/question-02.pdf}
\caption*{图2-1　Michelson 干涉仪的分束、往返反射和观察方向}
\end{figure}

\newpage

\section{第 3 题}\label{ux7b2c-3-ux9898}

半径为 \(r\)
的圆柱被两条不可伸长的绳沿两条切线拉动，两条绳的切线方向夹角为
\(\alpha\)，从切点到牵引端的直线段长度分别为
\(l_1\)、\(l_2\)。牵引端沿各自绳方向的速度为 \(v_1\)、\(v_2\)。求：

（1）圆柱不自转时，圆心速度大小 \(v\) 和绳段 \(l_1\) 的转动角速度
\(\omega_1\)；

（2）圆柱以角速度 \(\omega\) 顺时针转动时，圆心速度大小 \(v_O\) 和绳段
\(l_1\) 的转动角速度 \(\omega_1'\)。

\begin{figure}[H]
\centering
\includegraphics[width=0.468\linewidth,height=\textheight,keepaspectratio,alt={图3-1　圆柱上的两条切绳、牵引速度和几何夹角}]{../figures/question-03.pdf}
\caption*{图3-1　圆柱上的两条切绳、牵引速度和几何夹角}
\end{figure}

\newpage

\section{第 4 题}\label{ux7b2c-4-ux9898}

长为 \(l\)、质量为 \(m\)
的均匀细杆下端与地面铰接，初始竖直静止，受微小扰动后向下倒。当杆与竖直方向夹角为
\(\theta\) 时，取距铰点为 \(x\) 的截面，求：

（1）该截面内力的径向分量 \(F_{x1}\)，规定拉力为正；

（2）该截面内力的切向分量大小 \(F_{x2}\)；

（3）若结构耐剪切但不耐拉伸，求最容易发生拉断的位置 \(x_0/l\)。

\begin{figure}[H]
\centering
\includegraphics[width=0.58\linewidth,height=\textheight,keepaspectratio,alt={图4-1　倒杆角度、截面位置及内力分量定义}]{../figures/question-04.pdf}
\caption*{图4-1　倒杆角度、截面位置及内力分量定义}
\end{figure}

\newpage

\section{第 5 题}\label{ux7b2c-5-ux9898}

一焦距为 \(f\) 的凸透镜与一焦距为 \(-f\) 的凹透镜共轴放置，两透镜间距为
\(f\)。轴上物点位于凸透镜左侧，求：

（1）物距为 \(1.5f\) 时，最终像的位置；

（2）物距为 \(2f\) 时，出射光的状态；

（3）若最终像位于凹透镜右侧、且像到凹透镜的距离等于物体到凸透镜的距离，求该距离；

（4）在凹透镜右侧放置一块与透镜垂直的平面镜，使光线沿原路返回。物体放在何处时，返回光束仍在原物位置成像？

\begin{figure}[H]
\centering
\includegraphics[width=0.8\linewidth,height=\textheight,keepaspectratio,alt={图5-1　共轴正、负透镜组及平面镜的配置}]{../figures/question-05.pdf}
\caption*{图5-1　共轴正、负透镜组及平面镜的配置}
\end{figure}

\newpage

\section{第 6 题}\label{ux7b2c-6-ux9898}

半径为 \(R\) 的均匀绝缘球带总电荷
\(-Q\)。沿一条半径开有很窄的光滑隧道，隧道对原电荷分布的影响忽略不计。一只质量为
\(m\)、带正电荷 \(q\) 的小物体可在隧道内运动，重力加速度为
\(g\)，静电力常量为 \(k\)。求：

（1）小物体从球的最低点出发，恰能到达球心且速度为零时的 \(q\)；

（2）在第（1）问条件下，小物体从最低点由静止释放后首次到达球心所需时间；

（3）忽略重力，令小物体的正电荷量为 \(Q\)，从球心正下方距球心 \(2R\)
处由静止释放，求它到达球心时的速度；

（4）第（3）问中，小物体从 \(2R\) 处运动到球面的时间；

（5）第（3）问中，小物体从球面运动到球心的时间。

\begin{figure}[H]
\centering
\includegraphics[width=0.72\linewidth,height=\textheight,keepaspectratio,alt={图6-1　带电绝缘球、径向隧道与小物体的初始位置}]{../figures/question-06.pdf}
\caption*{图6-1　带电绝缘球、径向隧道与小物体的初始位置}
\end{figure}

\newpage

\section{第 7 题}\label{ux7b2c-7-ux9898}

两根完全相同的均匀细杆 \(A\)、\(B\)，质量均为 \(m\)，长度均为
\(l\)。求：

（1）细杆绕质心及绕端点的转动惯量；

（2）两杆均在竖直平面内由水平位置静止释放。杆 \(A\) 的一端固定铰接，杆
\(B\)
的相应端点被约束在同一水平直线上作无摩擦运动。当两杆分别转到与竖直方向夹角
\(\theta\) 时，求角速度之比 \(\omega_A/\omega_B\)；

（3）两杆在光滑水平面上首尾铰接并处于同一直线，杆 \(A\)
的左端固定铰接。在杆 \(B\) 上距其右端 \(h\) 处施加垂直于杆的瞬时冲量
\(I\)。若冲量后两杆仍保持共线转动，求 \(h\)及共同角速度\(\omega\)。

\begin{figure}[H]
\centering
\includegraphics[width=0.86\linewidth,height=\textheight,keepaspectratio,alt={图7-1　单杆转动、滑动端约束和双杆冲量模型}]{../figures/question-07.pdf}
\caption*{图7-1　单杆转动、滑动端约束和双杆冲量模型}
\end{figure}

\newpage

\section{第 8 题}\label{ux7b2c-8-ux9898}

在惯性系 \(S\) 中，\(t=0\) 时三个质点的位置分别为

\[
A(0,0),\qquad B(l,0),\qquad C(l,l),
\]

其中 \(l=10\ \mathrm{ly}\)。质点 \(A\) 以速度 \(v=0.6c\) 沿 \(+x\)
方向运动，质点 \(B\) 以相同速率沿 \(-x\) 方向运动，质点 \(C\)
以相同速率沿 \(-y\) 方向运动。\(t=0\) 时 \(A\) 发出一束光，使其恰好与
\(C\) 相遇。求：

（1）光在 \(S\) 系中的发射方向；

（2）光与 \(C\) 相遇的时间；

（3）在随 \(A\) 运动的惯性系 \(S'\) 中，光的发射方向及相遇时间；

（4）在 \(S'\) 系中 \(t'=0\) 时，\(A\) 与 \(C\) 的距离。

\begin{figure}[H]
\centering
\includegraphics[width=0.524\linewidth,height=\textheight,keepaspectratio,alt={图8-1　三质点初始位置、运动方向与追光几何}]{../figures/question-08.pdf}
\caption*{图8-1　三质点初始位置、运动方向与追光几何}
\end{figure}
\end{document}
